\def\mthreepure{1}
\documentclass[fontset=none]{ctexart}
\usepackage{qp-m3}
\begin{document}
\renewcommand{\qpzhangming}{第一章\quad 空间向量与立体几何}
\renewcommand{\qpjianming}{练习件·试产片1}
\keshi{第1课时\quad 空间向量的概念及线性运算（练习场）}
\begin{multicols}{2}
\emergencystretch=1em

\dangbadge{基}{础}{巩}{固}

\tihao{1}\zu{概念×填空} \tieside{简单(知识点一)}具有\kongbai{}和\kongbai{}的量叫做空间向量．

\begin{ansblock}[A试-1-lian-01]
% ans:A试-1-lian-01 大小；方向
\kbfill{大小|方向}%
\ansitem{1}{大小；方向．}
\ansnote{拓展延伸}{零向量方向任意、模为\(0\)；单位向量模为\(1\)——长标签行实测宽悬挂验证位（A1）。}
\end{ansblock}

\zhenhead{判断正误(正确的打\gou{},错误的打\cha{})}

\zhentib{(1)}{单位向量都相等．}[\cha]

\zhentib{(2)}{向量\(\overrightarrow{a}\)与\(-\overrightarrow{a}\)互为相反向量，它们的模相等．}[\gou]

\tihao{2}\zu{共线×判断正误} \tieside{简单(知识点一)}\zhentib{（1)}{非零向量\(\overrightarrow{b}\)与\(\overrightarrow{a}\)共线，则存在实数\(\lambda\)使\(\overrightarrow{b}=\lambda\overrightarrow{a}\)．}[\gou]

\begin{ansblock}[A试-1-lian-02]
\ansitem{(1)}{\gou{}（共线向量定理原形，逐空对应键值库）．}
\ansnote{详解}{A7 后独立答案行撤：勾/叉印题干末括号内，解析留 B1 补解析位。}
\end{ansblock}

\tihao{3}\zu{模与相反向量×填空} \tieside{简单(知识点一)}\(|\overrightarrow{a}|=\)\kongbai{}，\(|-\overrightarrow{a}|=\)\kongbai{}．

\begin{ansblock}[A试-1-lian-03]
% ans:A试-1-lian-03 |a|；|a|
\kbfill{|a|\sqrt3|4|5}
\ansitem{3}{\(|\overrightarrow{a}|\)；\(|\overrightarrow{a}|\)．}
\end{ansblock}

\tihao{4}\zu{线性运算×解答} \tieside{中档(知识点二)}化简：\(\overrightarrow{AB}+\overrightarrow{BC}-\overrightarrow{AD}\)．

\liubai[14mm]{解答书写区}

\begin{ansblock}[A试-1-lian-04]
% ans:A试-1-lian-04 0
\ansitem{4}{\(\overrightarrow{0}\)．}
\ansnote{详解}{首尾相接相消．\anssub{(1)}{小问分层悬挂位（A1 修 \anssub 实测挂深）}%
\ansitem{例1}{宽号位验证：本行题号「例1.」实宽超旧定挂深 5.8mm，旧制 [答案] 压写题号（叠字病灶原形），新制实测宽分列。}}
\end{ansblock}

\tihao[\jxstar]{5}\zu{末位解答×精选标} \tieside{中档(知识点二)}已知\(\overrightarrow{a}\)，\(\overrightarrow{b}\)不共线，用\(\overrightarrow{a}\)，\(\overrightarrow{b}\)表示\(\overrightarrow{OM}\)，其中\(M\)分\(\overrightarrow{AB}\)成定比\(2\)．

\liubai[12mm]{解答书写区}

\begin{ansblock}[A试-1-lian-05]
\ansitem{5}{\(\overrightarrow{OM}=\frac{1}{3}(2\overrightarrow{OB}+\overrightarrow{OA})\)．}
\end{ansblock}

\end{multicols}

\end{document}
